\begin{definition}[Lamprecht formula with stationary classes]\label{mod:lamprecht-formula}
Let $F$ be a nonarchimedean local field.  Let $\theta$ and $\psi$ have
exact multiplicative and additive conductors, and write
\[
 q=m_F(\theta)=2d+\varepsilon>1,
 \qquad \varepsilon\leq 1,
 \qquad r=d+\varepsilon.
\]
Thus $\varepsilon\in\{0,1\}$ and $r$ is the minimal stationary depth.
Fix an admissible $\Gamma\in F^\times$, so
$\operatorname{ord}_F(\Gamma)=q+n_F(\psi)$.  The stationary numerator used
throughout the construction is the quotient class
\[
 \operatorname{St}_{\Gamma,r}(\theta)
   \in \mathcal O_F/\mathfrak p_F^d.
\]
A stationary representative is a deliberately supplied
$\beta\in\mathcal O_F$ together with a proof that its class is the displayed
stationary class; no declaration chooses such a representative.  Lean proves
that every permitted representative has exact order zero and hence determines
a local unit.  Put
\[
 E_\beta=\psi(\beta/\Gamma)\theta(\beta)^{-1}.
\]

If $\varepsilon=0$, Lean evaluates the unnormalised finite Gauss sum and its
phase exactly:
\[
 G_\Gamma(\theta,\psi)=q_F^dE_\beta,
 \qquad
 \Delta_F^{\mathrm{fin}}(\theta,\psi;\Gamma)
   =\theta(\Gamma)E_\beta.
\]
Here $q_F=|k_F|$; it is distinct from the conductor $q$, and the positive
factor $q_F^d$ occurs in the raw sum before disappearing under the phase.

Suppose $\varepsilon=1$.  Let $\delta\in F^\times$ be any element satisfying
$\operatorname{ord}_F(\delta)=d$.  Multiplication by $\delta^2$ descends from
arbitrary integral lifts to the critical quotient.  Pairing this map with the
stationary quotient class defines a representative-independent, nontrivial
additive character $\psi_0$ of $k_F$ such that, for every integral lift
$\widetilde z$ and every permitted $\beta$,
\[
 \psi_0(\bar z)=
 \psi(\beta\delta^2\widetilde z/\Gamma).
\]
Lean constructs the genuine Hasse function
\[
 \phi_\beta(\bar x)=
 \psi(\delta\beta\widetilde x/\Gamma)
 \theta(1+\delta\widetilde x)^{-1}
\]
and proves this value for every integral lift.  It is nowhere zero and has the
manuscript's positive polar law
\[
 \phi_\beta(x+y)=
 \phi_\beta(x)\phi_\beta(y)\psi_0(xy).
\]
Finite additive-character orthogonality gives
\[
 \left\lVert\sum_{x\in k_F}\phi_\beta(x)\right\rVert=\sqrt{q_F},
 \qquad
 \sum_{x\in k_F}\phi_\beta(x)\ne0.
\]
The exact raw and phase formulas are
\[
 G_\Gamma(\theta,\psi)
   =q_F^dE_\beta\sum_{x\in k_F}\phi_\beta(x),
\]
\[
 \Delta_F^{\mathrm{fin}}(\theta,\psi;\Gamma)
   =\theta(\Gamma)E_\beta
      \Aphase\!\left(\sum_{x\in k_F}\phi_\beta(x)\right).
\]

Finally let $\beta'=\beta+h$ be another representative of the same stationary
class.  In the even case $E_{\beta'}=E_\beta$.  In the odd case Lean constructs
a residue class $a\in k_F$ by reducing the proved-integral quotient
$h/(\beta\delta)$.  The residual polar character is unchanged, while Lean
proves
\[
 \phi_{\beta'}=T_a\phi_\beta,
 \qquad
 \Aphase\!\left(\sum\phi_{\beta'}\right)
   =\phi_\beta(a)^{-1}
      \Aphase\!\left(\sum\phi_\beta\right),
 \qquad
 E_{\beta'}=\phi_\beta(a)E_\beta.
\]
Consequently the complete elementary--Hasse product is independent of the
chosen representative.  The two odd factors are transported together; Lean
does not assert that either factor is separately representative-independent.
\LeanFile{LanglandsFirstMainLemma/Lamprecht/Formula.lean}
\lean{LanglandsFirstMainLemma.lamprechtFormula_stationaryDepth}
\lean{LanglandsFirstMainLemma.StationaryParameter}
\lean{LanglandsFirstMainLemma.lamprechtStationaryRepresentativeUnit_ord}
\lean{LanglandsFirstMainLemma.lamprechtResidualVariable}
\lean{LanglandsFirstMainLemma.lamprechtResidualVariable_integral_lift}
\lean{LanglandsFirstMainLemma.lamprechtResidualAddChar}
\lean{LanglandsFirstMainLemma.lamprechtResidualAddChar_integral_lift}
\lean{LanglandsFirstMainLemma.lamprechtResidualAddChar_ne_one}
\lean{LanglandsFirstMainLemma.lamprechtOdd_residualCoefficient_ord}
\lean{LanglandsFirstMainLemma.lamprechtHasseFunction}
\lean{LanglandsFirstMainLemma.lamprechtHasseFunction_integral_lift}
\lean{LanglandsFirstMainLemma.lamprechtElementaryFactor}
\lean{LanglandsFirstMainLemma.lamprechtEven_elementary_representative_independent}
\lean{LanglandsFirstMainLemma.lamprechtOdd_representative_transport}
\lean{LanglandsFirstMainLemma.lamprechtOdd_hasseSum_ne_zero}
\lean{LanglandsFirstMainLemma.lamprechtOdd_hasseSum_norm}
\lean{LanglandsFirstMainLemma.lamprechtOdd_hasseSumPhase_translate}
\lean{LanglandsFirstMainLemma.lamprechtOdd_complete_representative_independent}
\lean{LanglandsFirstMainLemma.finiteGaussSum_lamprecht_even}
\lean{LanglandsFirstMainLemma.finiteGaussSum_lamprecht_odd}
\lean{LanglandsFirstMainLemma.lamprechtEven}
\lean{LanglandsFirstMainLemma.lamprechtOdd}
\leanok
\SourceText{Lemma~\ref{lem:stationary-class-existence};
Lemma~\ref{lem:stationary-class-calculus}; Theorem~\ref{thm:Lamprecht};
Lemma~\ref{lem:complete-stationary-invariant}.}
\uses{mod:lamprecht-stationaryclasscalculus,mod:delta-finitedefinition,mod:finitefield-quadraticphase,mod:finitefield-hassefunction}
\FormalizationContract The theorem concerns the representative-free,
unnormalised finite Gauss sum and the computational local constant
$\Delta_F^{\mathrm{fin}}$ from the preceding node.  Its stationary datum and
residual polar character are quotient-level constructions.  Every use of a
field-valued stationary numerator quantifies over a supplied representative
and proves its exact order.  The odd residual refinement is a genuine Hasse
function for a proved nontrivial additive character, with arbitrary-lift
evaluation, nonvanishing sum, and positive polar sign.  Under a change of
stationary representative, the elementary value and Hasse-sum phase are
transported and cancelled as a single complete invariant.
\end{definition}
